% TOP_PROBLEM: {"id":30007613,"problem_number":"LOCAL-30007613","title":"Exchange-rate disconnect","category_id":21,"category":{"id":21,"name":"economics","display_name":"Economics","description":"Open economic theory statements, published research agendas, and proposed scoped specifications; question provenance and historical priority are qualified per record.","slug":"economics","order_index":21,"created_at":"2026-10-08T00:00:00Z"},"status":"research_question","external_url":null,"legacy_ids":[],"published":false,"statement":"In the explicitly specified setting: What share of exchange-rate movements is driven by fundamentals versus financial-demand shocks and intermediary frictions? The mathematical statement above fixes the target, assumptions and scope of this catalogue formulation.","economics":{"original_id":102,"field":"International finance and exchange rates","kind":"identification","formalization_basis":"adapted_research_question","source_status":"explicit_open","priority_status":"earliest_found","priority_note":"The 2000 abstract names an umbrella puzzle whose empirical antecedents are earlier. Priority for the catalogue shock-share question is not established.","earliest_verified_reference":{"authors":["Maurice Obstfeld","Kenneth Rogoff"],"title":"The Six Major Puzzles in International Macroeconomics: Is There a Common Cause?","year":2000,"url":"https://www.journals.uchicago.edu/doi/10.1086/654423","doi":"10.1086/654423","locator":"Abstract, paragraph 2; article pp.339–390","evidence_summary":"Names the exchange-rate disconnect puzzle and its empirical antecedents."},"current_reference":{"authors":["Ryan Chahrour","Vito Cormun","Pierre De Leo","Pablo Guerrón-Quintana","Rosen Valchev"],"title":"Exchange Rate Disconnect Revisited","year":2025,"url":"https://chahrour.github.io/papers/Disconnect_revisited.pdf","doi":"10.3386/w32596","locator":"November 21, 2025 manuscript; abstract and introduction, pp.1–4","evidence_summary":"Noisy future-productivity news explains a large sample-specific fluctuation share."},"formalization_note":"Includes noisy news about future fundamentals because current evidence challenges a literal total disconnect. The proposed variance decomposition is identification-dependent. Independent math audit: the variance-share denominator is required to be positive and the horizon is positive."},"statusQualification":"Published question or research agenda verified at the cited locator. The mathematical specification is adapted for this catalogue and includes additional assumptions; source publication does not establish first-ever priority or certify this exact model remains unsolved. Includes noisy news about future fundamentals because current evidence challenges a literal total disconnect. The proposed variance decomposition is identification-dependent. Independent math audit: the variance-share denominator is required to be positive and the horizon is positive.","statusReviewedAt":"2026-10-08"}
% FIRST_POSTED: 2026-10-08
% ENTRY_KIND: statement_only
% !TeX program = lualatex
\documentclass[11pt]{article}
\usepackage{fontspec}
\usepackage[a4paper,margin=25mm]{geometry}
\usepackage{amsmath,amssymb,mathtools}
\usepackage{xurl}
\usepackage[hidelinks]{hyperref}
\newcommand{\sref}[2]{\hyperlink{source:#1:#2}{[S#2]}}
\setlength{\emergencystretch}{3em}
\begin{document}
% problemId:30007613
\section*{Exchange-rate disconnect}\label{30007613}
\subsection{Definitions and mathematical statement}
Fix exchange-rate data \(s_t\), current and future productivity, consumption, trade, financing conditions and measured expectations for specified countries and dates. Let \(y_t\) collect these observables. A structural representation is
\[y_t=\sum_{h\geq0}B_h\varepsilon_{t-h},\qquad E[\varepsilon_t\varepsilon_t']=I,\]
where \(B_h\) are response matrices and \(\varepsilon_t\) are mean-zero structural shocks divided into current fundamentals, noisy news about future fundamentals, financial-demand shocks and intermediary-capacity shocks. The summation must converge in mean square; \(I\) is the identity covariance matrix. The task is to identify the exchange-rate forecast-error variance shares attributable to each group at declared horizon \(H\geq1\), restricting attention to positive exchange-rate forecast-error variance,
\[\theta_g(H)=\frac{\sum_{h=0}^{H-1}\sum_{k\in g}(B_h)_{s,k}^2}{\sum_{h=0}^{H-1}\sum_k(B_h)_{s,k}^2}.\]
Here \((B_h)_{s,k}\) is the response of \(s\) to shock \(k\) and group \(g\) is specified before estimation. Provide restrictions distinguishing news from contemporaneous fundamentals and financial shocks; sign restrictions alone may yield sets rather than unique shares. Examine whether apparent weak contemporaneous correlations reflect different response horizons rather than absence of fundamental links. Report identified sets and robustness to omitted shocks, exchange regimes and survey measurement errors. The decomposition requires an identification design and does not follow from low exchange-rate correlations alone.

\par\textbf{Formulation and scope.} Includes noisy news about future fundamentals because current evidence challenges a literal total disconnect. The proposed variance decomposition is identification-dependent. Independent math audit: the variance-share denominator is required to be positive and the horizon is positive.

\par\textbf{Open-question provenance.} An explicit question or research agenda was verified at the source locator. This does not by itself certify that every later version remains unresolved \sref{30007613}{1}.

\par\textbf{Historical priority.} The 2000 abstract names an umbrella puzzle whose empirical antecedents are earlier. Priority for the catalogue shock-share question is not established.

\subsection{Short English statement}
In the explicitly specified setting: What share of exchange-rate movements is driven by fundamentals versus financial-demand shocks and intermediary frictions? The mathematical statement above fixes the target, assumptions and scope of this catalogue formulation.

\subsection{Sources}
\begin{itemize}\raggedright
\item[\textnormal{[S1]}]\hypertarget{source:30007613:1}{} Maurice Obstfeld, Kenneth Rogoff. 2000. The Six Major Puzzles in International Macroeconomics: Is There a Common Cause?. Abstract, paragraph 2; article pp.339–390.
\url{https://www.journals.uchicago.edu/doi/10.1086/654423}.
DOI: 10.1086/654423.
{\small\textit{Earliest verified question/agenda reference; historical priority qualified below. Names the exchange-rate disconnect puzzle and its empirical antecedents.}}

\item[\textnormal{[S2]}]\hypertarget{source:30007613:2}{} Ryan Chahrour, Vito Cormun, Pierre De Leo, Pablo Guerrón-Quintana, Rosen Valchev. 2025. Exchange Rate Disconnect Revisited. November 21, 2025 manuscript; abstract and introduction, pp.1–4.
\url{https://chahrour.github.io/papers/Disconnect_revisited.pdf}.
DOI: 10.3386/w32596.
{\small\textit{Supporting or current literature; see evidence qualification. Noisy future-productivity news explains a large sample-specific fluctuation share.}}
\end{itemize}
\end{document}
